\begin{center}
\LARGE
Nonlinear network theory, differential geometric approach
                     to nonlinear networks
\end{center}

\begin{center}
\Large
Wolfgang Mathis

\end{center}

\noindent
Date:  July 19th (Friday)\\
Time: 10:30-11:00\\

\begin{center}
\large
Abstract
\end{center}

In order to design electronic circuits the application of rather heuristic methods are used by many
designers. Unfortunately these methods are not very successful in the case of circuits where the
nonlinearities are essential; oscillators and digital subcircuits are of this kind. Furthermore a
circuit simulator is used very often within a design process but it is only applicable of all circuit
parameters are fixed and in certain cases this tool runs into difficulties. 

In this talk we present some fundamentals of circuit theory based on modern differential geometry.
An advantage of this approach is that nongeneric cases can be classified. If a circuit simulator is
applied to a circuit that is "near" such an untypical case it can be run into the mentioned problems.
Further this geometric approach is suitable to construct a systematic design process for sinusoidal
oscillators. Some of the design algorithms are discussed that are implemented in
MATHEMATICA. 

